\chapter{Моментная реализация одной нити и минимальный запрет}
\label{ch:v6-single-thread-moment-realization-gate}

\section{Локальная мера проходов}

Пусть одна ориентированная глобальная нить многократно пересекает малую
пространственную ячейку. После усреднения её проходы задают
неотрицательные веса $w_a$ направлений $n_a$, где
\begin{equation}
 n_1=\frac{(1,1,1)}{\sqrt3},\quad
 n_2=\frac{(1,-1,-1)}{\sqrt3},\quad
 n_3=\frac{(-1,1,-1)}{\sqrt3},\quad
 n_4=\frac{(-1,-1,1)}{\sqrt3}.
\end{equation}
Они удовлетворяют
\begin{equation}
 \sum_a n_a=0,
 \qquad
 n_a\mathbin{\cdot}n_b=-\frac13\quad(a\ne b).
\end{equation}

Обычные первые три момента одной локальной плотности равны
\begin{align}
 J_i&=\sum_a w_a(n_a)_i,\\
 R_{ij}&=\sum_a w_a(n_a)_i(n_a)_j,\\
 M_{ijk}&=\sum_a w_a(n_a)_i(n_a)_j(n_a)_k.
\end{align}
Для связной цепи $J$, плотность и тензорный порядок не являются
независимыми полями: связность линии порождает векторные и тензорные
законы сохранения
\cite{single-thread-moment-svensek2013,single-thread-moment-svensek2016}.

\section{Точное воспроизведение второго момента}

Упорядоченное состояние проекта имеет спектр
\begin{equation}
 (a,b,b)=(0.9121666003,0.04391669984,0.04391669984).
\end{equation}
Выберем ось $n_1$ и потребуем циклически симметричные веса
$(p,q,q,q)$. Равенство
\begin{equation}
 R_*=a,n_1n_1^{\mathsf T}+b(I-n_1n_1^{\mathsf T})
\end{equation}
фиксирует
\begin{equation}
 p=\frac{9a-1}{8}=0.9011874254,
 \qquad
 q=\frac{1-p}{3}=0.03293752488.
\end{equation}
Остаток второго момента равен нулю с машинной точностью. Однако тот же
набор проходов несёт полярный ток
\begin{equation}
 \|J\|=0.8682499005.
\end{equation}
После бесследовой проекции его третий момент не совпадает с каноническим
\begin{equation}
 T_*=\sum_{a=1}^{4}n_a^{\otimes3}.
\end{equation}
Лучшее масштабное совмещение имеет косинус $0.789808034$ и относительный
остаток $0.613354114$.

\section{Нулевой ток фиксирует изотропию}

Матрица ограничений, состоящая из нормировки
$\sum_a w_a=1$ и трёх компонент $\sum_aw_an_a=0$, имеет ранг четыре.
Поэтому решение единственно:
\begin{equation}
 w_1=w_2=w_3=w_4=\frac14.
\end{equation}
Оно действительно даёт ненулевой канонический третий момент
\begin{equation}
 \left\|\sum_a\frac14n_a^{\otimes3}\right\|
 =0.4714045208,
\end{equation}
но второй момент становится точно изотропным:
\begin{equation}
 \sum_a\frac14 n_an_a^{\mathsf T}=\frac13I,
 \qquad Q=0.
\end{equation}

Если вместо ориентированных проходов взять Real-сбалансированные пары
$\{n_a,-n_a\}$, то одновременно исчезают первый и все нечётные обычные
моменты. Численно $\|M_3\|=1.70\cdot10^{-17}$, поэтому такая пара не
воспроизводит $T_*$.

\begin{theorem}[Запрет общей сырой моментной плотности]
Одно неотрицательное распределение на четырёх тетраэдрических направлениях
не может одновременно воспроизвести одноосный второй момент $Q_*$,
канонический тетраэдрический третий момент $T_*$ и нулевой полярный ток.
Условие $J=0$ фиксирует равные веса и тем самым обращает $Q$ в нуль;
Real-баланс направлений дополнительно обращает в нуль $T$.
\end{theorem}

\section{Что именно закрыто}

Результат не опровергает одну глобальную нить. Он опровергает только
минимальное отождествление
\begin{equation}
 Q=\text{сырой второй момент},
 \qquad
 T=\text{сырой третий момент}
\end{equation}
одной и той же локальной плотности проходов без дополнительных меток.

Минимальное допустимое продолжение состоит в том, что $Q$ измеряет
плотность направлений, а $T$ использует на тех же проходах выведенный
фазовый, ориентационный или связностный вес. Такой вес не может быть
произвольным: он должен следовать из уже существующих $Z_3$-голономии,
Real-градуировки или локального правила сшивки.

\begin{nogocriterion}
Запрещено спасать моментную реализацию независимыми весами для $Q$ и $T$,
если их отношение просто подбирается. Следующий гейт проходит только при
каноническом весе из существующей фазы или связности.
\end{nogocriterion}

\section{Следующий тест}

Проверить фазово-взвешенный третий момент на тех же четырёх проходах.
Нулевая гипотеза: характер остаточной $Z_3$-голономии или Real-пара
однозначно задаёт веса, которые сохраняют $R_*$, обнуляют физический
полярный ток и воспроизводят направление $T_*$ без нового параметра.

Стоп-критерий: если все канонические характеры либо зануляют $T$, либо
оставляют независимый модуль весов, прямой мост одной нити к существующему
$Q+T$-вакууму закрывается.

\begin{protocol}
\begin{enumerate}
 \item упорядоченный второй момент: \textbf{воспроизведён точно};
 \item требуемые веса: \textbf{$(0.9011874,0.0329375^3)$};
 \item полярный ток этих весов: \textbf{$0.8682499$};
 \item совпадение сырого третьего момента с $T_*$: \textbf{нет};
 \item единственное решение с нулевым током: \textbf{равные веса};
 \item $Q$ при равных весах: \textbf{ноль};
 \item $T$ у Real-сбалансированной пары: \textbf{ноль};
 \item простая общая моментная плотность: \textbf{закрыта};
 \item одна глобальная нить: \textbf{не опровергнута};
 \item минимальное продолжение: \textbf{фазово-взвешенный момент};
 \item рождение материи: \textbf{не закрыто}.
\end{enumerate}
\end{protocol}

Машинный сертификат сохранён в
\path{s2t_v6_single_thread_moment_realization_gate_results.json}.

\begin{thebibliography}{9}
\bibitem{single-thread-moment-svensek2013}
D.~Sven\v{s}ek, G.~M.~Grason, R.~Podgornik,
\emph{Tensorial conservation law for nematic polymers},
arXiv:1309.7125.

\bibitem{single-thread-moment-svensek2016}
D.~Sven\v{s}ek, R.~Podgornik,
\emph{A generalized conservation law for main-chain polymer nematics},
arXiv:1601.06582.

\bibitem{single-thread-moment-radzihovsky2000}
L.~Radzihovsky, T.~C.~Lubensky,
\emph{Fluctuation-Driven First-Order Isotropic-to-Tetrahedratic Phase Transition},
arXiv:cond-mat/0008325.

\bibitem{single-thread-moment-lubensky2002}
T.~C.~Lubensky, L.~Radzihovsky,
\emph{Theory of Banana Liquid Crystal Phases and Phase Transitions},
arXiv:cond-mat/0205171.
\end{thebibliography}